\section{总结与延伸}

本讲完整串起了三条路线：用 learned model 扩充数据、用任务条件化构建 generalist policy、用 hindsight relabeling 把失败轨迹转化为有效监督。三者虽然表面不同，本质都在回答同一个问题：\textbf{如何让有限经验服务更多决策。}

\subsection{本讲核心总结表}

\begin{table}[H]
\centering
\begin{tabular}{p{0.18\textwidth} p{0.22\textwidth} p{0.22\textwidth} p{0.24\textwidth}}
\toprule
主题 & 核心问题 & 关键技巧 & 主要限制 \\
\midrule
Model-based RL & 如何降低真实交互成本 & learned model + short synthetic rollout & 模型偏差、长 horizon 误差累积 \\
多任务 RL & 如何在任务间共享经验与技能 & task conditioning + shared policy & 负迁移、任务分布不均衡 \\
Goal-conditioned RL & 如何学习通用到达策略 & 目标作为条件输入策略 & 目标表示和奖励设计困难 \\
HER & 如何让失败轨迹也有训练价值 & hindsight relabeling + off-policy reuse & 需要可重算奖励和共享动力学 \\
\bottomrule
\end{tabular}
\caption{Lecture 12 的方法脉络}
\end{table}

\subsection{进一步思考}

\begin{itemize}
    \item 当 task descriptor 从 one-hot 变成自然语言后，多任务 RL 与 foundation model 的界线正在模糊。
    \item HER 的成功提醒我们：很多看似 ``失败'' 的数据，其实只是缺少恰当的标签和解释方式。
    \item 课程最后预告的下一个问题非常自然：既然可以在多个任务间共享知识，那么能否进一步做到对全新任务的快速适应？
\end{itemize}

\subsection{从多任务共享走向快速适应}

这节课结尾其实埋下了一个很重要的研究转折：如果多任务 RL 解决的是 ``在一批已知任务之间共享结构''，那么下一步自然就会问，\textbf{面对一个此前从未见过的新任务，能不能像 in-context learning 那样快速适配，而不是重新训练整个策略？}

\begin{figure}[H]
\centering
\includegraphics[width=0.9\textwidth]{slides-images/slide-29.jpg}
\caption{课程结尾抛出的下一个问题：能否对新任务实现快速适应}
\end{figure}

从这个角度看，本讲中的三条主线恰好构成了后续研究的基础：
\begin{enumerate}
    \item Model-based RL 让 agent 更善于复用已有经验和环境结构。
    \item 多任务 RL 让策略学会在多个任务间共享参数与行为模式。
    \item HER 说明同一条轨迹可以在不同目标解释下反复提供监督。
\end{enumerate}

它们共同指向一个更大的问题：\textit{如果经验、任务条件和目标解释都能重用，那么策略是否也应当具备更强的 test-time adaptation 能力？}

\subsection{实现这讲方法时最容易踩的坑}

这节课虽然概念上分成三部分，但真正写代码时常常会把坑踩在数据管线上。尤其是 HER 和多任务条件化策略，往往不是算法公式错了，而是 replay buffer、task descriptor 和 reward recomputation 之间没有接好。

\begin{figure}[H]
\centering
\includegraphics[width=0.9\textwidth]{slides-images/slide-28.jpg}
\caption{Lecture 12 的主线回看：model-based augmentation、task conditioning 与 hindsight relabeling}
\end{figure}

\begin{table}[H]
\centering
\begin{tabular}{p{0.2\textwidth} p{0.28\textwidth} p{0.28\textwidth}}
\toprule
实现模块 & 最常见错误 & 更稳妥的做法 \\
\midrule
Synthetic rollout & rollout 太长导致模型偏差主导训练 & 从真实数据中的中间状态出发，只做短 rollout \\
Task conditioning & task ID 和状态没有同步送入 policy / critic & 统一把 $(\bar{s}, z)$ 当作条件输入接口 \\
HER relabeling & 只改目标，不重算整条轨迹奖励 & 明确实现 reward recomputation 与 buffer 写回逻辑 \\
Replay buffer & 原始样本和 relabeled 样本分布失衡 & 记录采样比例，必要时做分层采样 \\
\bottomrule
\end{tabular}
\caption{Lecture 12 对应的实现清单}
\end{table}

\begin{knowledgebox}{拓展阅读}
\begin{itemize}
    \item Andrychowicz et al., ``Hindsight Experience Replay,'' NeurIPS 2017
    \item Schaul et al., ``Universal Value Function Approximators,'' ICML 2015
    \item Yu et al., ``Meta-World: A Benchmark and Evaluation for Multi-Task and Meta Reinforcement Learning,'' CoRL 2020
    \item Ghosh et al., ``Learning to Reach Goals via Iterated Supervised Learning,'' ICLR 2021
\end{itemize}
\end{knowledgebox}

\end{document}
